% GENERATED FILE. Edit canonical lessons, then run npm run generate:books.
% canonical-source-hash: fnv1a64:58ef0d4b052ec768
% canonical-lessons: LA-C473-rus, LA-C473-fundus, LA-C473-bubile, LA-C473-casa, LA-C473-tugurium

\chapter{Countryside, Farm, Ox stall, Cottage, Hut}
\label{ch:la-countryside-farm-ox-stall-cottage-hut}
\hlchaptermodality{473}

\begin{chapteropening}
\textbf{By the end of this chapter:} \emph{I can talk about the countryside and the buildings of a farm: the country, a farm, the ox stall, a cottage and a hut.}

Complete the last lesson of chapter 473: I can talk about the countryside and the buildings of a farm: the country, a farm, the ox stall, a cottage and a hut.
\end{chapteropening}

\section[rūs, rūris]{rūs, rūris — the countryside, the country; a farm}

\label{lesson:LA-C473-rus}



\begin{quote}
Before the new one: say the Latin for to spin, then the Latin for wool-working.
\end{quote}

\subsection*{You'll want to know: rūs, rūris}
\textbf{rūs, rūris} — \textquotedblleft{}the countryside, the country; a farm\textquotedblright{}.

\begin{quote}
Aestāte multī cīvēs rūs eunt.
\end{quote}

\begin{quote}
\emph{In summer many citizens go to the country.}
\end{quote}

Third declension, neuter: \emph{rūs, rūris n.} Like \emph{domus}, it takes no preposition: \emph{rūs} 'to the country', \emph{rūre} 'from the country', \emph{rūrī} 'in the country'. English \emph{rural} and \emph{rustic} come from it.

\begin{cousinweb}
\emph{rūs} is inherited from Proto-Italic \emph{rowos}, from PIE \emph{rewh\textsubscript{1}-} 'open space, field' (EtymDB). English \emph{rural} and, through \emph{rūsticus}, \emph{rustic} come from it; German \emph{Raum} 'space' is a cousin.
\end{cousinweb}

\begin{grammarlens}[title={What you've built}]
One more word for telling, explaining and describing at B1.
\end{grammarlens}

\subsection*{Your turn}
\begin{itemize}
\raggedright
  \item \emph{Say it:} \emph{rūs, rūris}
  \item \emph{Say it:} \emph{rūs, rūris}, once more
  \item \emph{Say it:} a sentence of your own with \emph{rūs, rūris}
  \item \emph{Recall:} say the Latin for to spin, then the Latin for wool-working, then say \emph{rūs, rūris} again
\end{itemize}

\begin{tcolorbox}[breakable,colback=teal!4,colframe=teal!35!black,title={Before you move on}]
What does rūs, rūris mean? (\textquotedblleft{}The countryside, the country; a farm\textquotedblright{}.) Say it once more.
\end{tcolorbox}
\section[fundus, fundī]{fundus, fundī — a farm, an estate; the bottom}

\label{lesson:LA-C473-fundus}



\begin{quote}
Before the new one: say the Latin for wool-working, then the Latin for the countryside.
\end{quote}

\subsection*{You'll want to know: fundus, fundī}
\textbf{fundus, fundī} — \textquotedblleft{}a farm, an estate; the bottom\textquotedblright{}.

\begin{quote}
Fundus noster parvus est sed bonus.
\end{quote}

\begin{quote}
\emph{Our farm is small but good.}
\end{quote}

Second declension, masculine: \emph{fundus, fundī m.} First the bottom of a thing, then the ground itself, a piece of land. English \emph{fund} and \emph{foundation} go back to the same word.

\begin{cousinweb}
\emph{fundus} is inherited from PIE \emph{b\textsuperscript{h}ud\textsuperscript{h}m\'{\={e}}n} 'bottom' (EtymDB), with Greek \emph{pythmēn} 'bottom' as a cousin. English \emph{fund} and \emph{profound} come from it.
\end{cousinweb}

\begin{grammarlens}[title={What you've built}]
One more word for telling, explaining and describing at B1.
\end{grammarlens}

\subsection*{Your turn}
\begin{itemize}
\raggedright
  \item \emph{Say it:} \emph{fundus, fundī}
  \item \emph{Say it:} \emph{fundus, fundī}, once more
  \item \emph{Say it:} a sentence of your own with \emph{fundus, fundī}
  \item \emph{Recall:} say the Latin for wool-working, then the Latin for the countryside, then say \emph{fundus, fundī} again
\end{itemize}

\begin{tcolorbox}[breakable,colback=teal!4,colframe=teal!35!black,title={Before you move on}]
What does fundus, fundī mean? (\textquotedblleft{}A farm, an estate; the bottom\textquotedblright{}.) Say it once more.
\end{tcolorbox}
\section[būbīle, būbīlis]{būbīle, būbīlis — an ox stall, a cowshed}

\label{lesson:LA-C473-bubile}



\begin{quote}
Before the new one: say the Latin for the countryside, then the Latin for a farm.
\end{quote}

\subsection*{You'll want to know: būbīle, būbīlis}
\textbf{būbīle, būbīlis} — \textquotedblleft{}an ox stall, a cowshed\textquotedblright{}.

\begin{quote}
Bovēs hieme in būbīlī manent.
\end{quote}

\begin{quote}
\emph{In winter the oxen stay in the cowshed.}
\end{quote}

Third declension, neuter: \emph{būbīle, būbīlis n.}, from \emph{bōs, bovis} 'ox'. A word of the farm writers.

\begin{cousinweb}
\emph{būbīle} is built on \emph{bōs, bovis} 'ox' with the suffix \emph{-īle} 'place for animals', as in \emph{ovīle} 'sheepfold' (Lewis \& Short).
\end{cousinweb}

\begin{grammarlens}[title={What you've built}]
One more word for telling, explaining and describing at B1.
\end{grammarlens}

\subsection*{Your turn}
\begin{itemize}
\raggedright
  \item \emph{Say it:} \emph{būbīle, būbīlis}
  \item \emph{Say it:} \emph{būbīle, būbīlis}, once more
  \item \emph{Say it:} a sentence of your own with \emph{būbīle, būbīlis}
  \item \emph{Recall:} say the Latin for the countryside, then the Latin for a farm, then say \emph{būbīle, būbīlis} again
\end{itemize}

\begin{tcolorbox}[breakable,colback=teal!4,colframe=teal!35!black,title={Before you move on}]
What does būbīle, būbīlis mean? (\textquotedblleft{}An ox stall, a cowshed\textquotedblright{}.) Say it once more.
\end{tcolorbox}
\section[casa, casae]{casa, casae — a cottage, a small house}

\label{lesson:LA-C473-casa}



\begin{quote}
Before the new one: say the Latin for a farm, then the Latin for an ox stall.
\end{quote}

\subsection*{You'll want to know: casa, casae}
\textbf{casa, casae} — \textquotedblleft{}a cottage, a small house\textquotedblright{}.

\begin{quote}
Pastor in casā parvā prope silvam habitat.
\end{quote}

\begin{quote}
\emph{The shepherd lives in a small cottage near the wood.}
\end{quote}

First declension, feminine: \emph{casa, casae f.} In Italian and Spanish it became the everyday word for 'house'.

\begin{cousinweb}
The origin of \emph{casa} is not settled: Lewis \& Short link it with a root meaning 'to cover', but EtymDB gives no parent. It became \emph{casa} 'house' in Italian, Spanish and Portuguese.
\end{cousinweb}

\begin{grammarlens}[title={What you've built}]
One more word for telling, explaining and describing at B1.
\end{grammarlens}

\subsection*{Your turn}
\begin{itemize}
\raggedright
  \item \emph{Say it:} \emph{casa, casae}
  \item \emph{Say it:} \emph{casa, casae}, once more
  \item \emph{Say it:} a sentence of your own with \emph{casa, casae}
  \item \emph{Recall:} say the Latin for a farm, then the Latin for an ox stall, then say \emph{casa, casae} again
\end{itemize}

\begin{tcolorbox}[breakable,colback=teal!4,colframe=teal!35!black,title={Before you move on}]
What does casa, casae mean? (\textquotedblleft{}A cottage, a small house\textquotedblright{}.) Say it once more.
\end{tcolorbox}
\section[tugurium, tuguriī]{tugurium, tuguriī — a hut, a shed}

\label{lesson:LA-C473-tugurium}



\begin{quote}
Before the new one: say the Latin for an ox stall, then the Latin for a cottage.
\end{quote}

\subsection*{You'll want to know: tugurium, tuguriī}
\textbf{tugurium, tuguriī} — \textquotedblleft{}a hut, a shed\textquotedblright{}.

\begin{quote}
Pauper agricola in tuguriō cum familiā vīvit.
\end{quote}

\begin{quote}
\emph{The poor farmer lives in a hut with his family.}
\end{quote}

Second declension, neuter: \emph{tugurium, tuguriī n.} A rough hut, poorer than a \emph{casa}.

\begin{cousinweb}
\emph{tugurium} is usually linked with \emph{tegō} 'to cover' (Lewis \& Short), so a shelter with a roof; EtymDB traces it to PIE \emph{(s)teg-} 'to cover with a roof', with Greek \emph{stegō} 'to cover' as a cousin.
\end{cousinweb}

\begin{grammarlens}[title={What you've built}]
One more word for telling, explaining and describing at B1.
\end{grammarlens}

\subsection*{Your turn}
\begin{itemize}
\raggedright
  \item \emph{Say it:} \emph{tugurium, tuguriī}
  \item \emph{Say it:} \emph{tugurium, tuguriī}, once more
  \item \emph{Say it:} a sentence of your own with \emph{tugurium, tuguriī}
  \item \emph{Recall:} say the Latin for an ox stall, then the Latin for a cottage, then say \emph{tugurium, tuguriī} again
\end{itemize}

\begin{tcolorbox}[breakable,colback=teal!4,colframe=teal!35!black,title={Before you move on}]
What does tugurium, tuguriī mean? (\textquotedblleft{}A hut, a shed\textquotedblright{}.) Say it once more.
\end{tcolorbox}
